\subsection{Second-order velocity structure functions and segmented power-law fits}
\label{sec:vsf2-methods}

We quantify the spatial scaling of the ionized-gas velocity field using the
second-order velocity structure function,
\begin{equation}
    S_2(\ell)
    =
    \left\langle
    \left| v(\mathbf{r}+\boldsymbol{\ell}) - v(\mathbf{r}) \right|^2
    \right\rangle ,
    \label{eq:vsf2}
\end{equation}
where $v(\mathbf{r})$ is the line-of-sight velocity at position
$\mathbf{r}$, $\ell = |\boldsymbol{\ell}|$ is the projected physical
separation, and the average is taken over all pairs of spaxels whose
separations fall in the same logarithmic radial bin. We compute $S_2$ for each
HII region and for each available emission-line tracer, using only spaxels that
pass the adopted signal-to-noise and spatial-selection criteria.

For each radial bin $j$, the measured structure function is
\begin{equation}
    S_{2,j}
    =
    \frac{1}{N_j}
    \sum_{i \in j}
    \left| \Delta v_i \right|^2 ,
    \label{eq:vsf2-binned}
\end{equation}
where $N_j$ is the number of velocity pairs in bin $j$ and
$\Delta v_i$ is the velocity difference for pair $i$. Because the pair values
are not statistically independent, we estimate the uncertainty in each radial
bin using a bootstrap over spaxels rather than over individual pairs. For each
bootstrap realization, we resample the input spaxels with replacement and
recompute $S_{2,j}$ in every radial bin. The adopted uncertainty is the
standard deviation of the bootstrap distribution,
\begin{equation}
    \sigma_{\mathrm{boot},j}
    =
    \mathrm{std}
    \left[
    S_{2,j}^{(b)}
    \right]_{b=1}^{N_{\mathrm{boot}}},
    \label{eq:vsf2-bootstrap}
\end{equation}
where $S_{2,j}^{(b)}$ is the value of the structure function in bin $j$ for
bootstrap realization $b$. This uncertainty captures both sampling variance
and the covariance introduced by reusing the same spaxels in many velocity
pairs.

The VSF scaling is fit in logarithmic space. We define
\begin{equation}
    x_j = \log_{10} \ell_j,
    \qquad
    y_j = \log_{10} S_{2,j},
    \qquad
    \sigma_{y,j}
    =
    \frac{\sigma_{\mathrm{boot},j}}{S_{2,j}\ln 10}.
    \label{eq:log-vsf}
\end{equation}
Only bins with finite positive $S_{2,j}$, finite positive uncertainty, and
sufficient pair counts are included in the fit.

To allow for possible turnover scales or changes in turbulent scaling, we fit
continuous segmented power laws. For a model with $m$ line segments and
$m-1$ breakpoints, the log-space model can be written as
\begin{equation}
    y(x)
    =
    a
    +
    b_1 x
    +
    \sum_{k=1}^{m-1}
    c_k \, \max(0, x-\tau_k),
    \label{eq:segmented-model}
\end{equation}
where $\tau_k$ are the breakpoint locations in $\log_{10}\ell$. This
parameterization enforces continuity at each breakpoint. The slope in the first
segment is $b_1$, while the slope after the $k$th breakpoint is
$b_1 + \sum_{q=1}^{k} c_q$.

For each VSF, we fit candidate models with one, two, and three line segments.
Breakpoints are searched over the observed radial bins, requiring a minimum
number of data points in every segment. For a fixed set of breakpoint
locations, the model is linear in the coefficients and is solved by weighted
least squares. The goodness of fit is
\begin{equation}
    \chi^2
    =
    \sum_j
    \left[
    \frac{y_j - y(x_j)}{\sigma_{y,j}}
    \right]^2 .
    \label{eq:chi2-vsf}
\end{equation}

We select the preferred number of segments using the Bayesian Information
Criterion (BIC),
\begin{equation}
    \mathrm{BIC}
    =
    \chi^2 + k \ln n,
    \label{eq:bic}
\end{equation}
where $n$ is the number of fitted radial bins and $k$ is the number of model
parameters. For a continuous segmented model with $m$ segments, we use
$k=2m$: one intercept, $m$ slope-equivalent parameters, and $m-1$ breakpoint
locations. Thus, the one-, two-, and three-segment models have
$k=2$, $4$, and $6$, respectively.

The model with the lowest BIC is favored, but an additional breakpoint is
accepted only when the BIC improvement over the best simpler model is
substantial. We define
\begin{equation}
    \Delta \mathrm{BIC}
    =
    \mathrm{BIC}_{\mathrm{simpler}}
    -
    \mathrm{BIC}_{\mathrm{complex}},
    \label{eq:delta-bic}
\end{equation}
and require $\Delta \mathrm{BIC} \geq 10$ for the more complex model to be
adopted. This threshold corresponds to strong evidence for the additional
breakpoint and guards against over-interpreting small improvements in
$\chi^2$.

For each accepted fit, we report the segment slopes, breakpoint locations, and
the BIC values for all tested models. Diagnostic plots show the measured
$S_2(\ell)$ values with their uncertainties together with the selected
segmented power-law model. These plots are used to verify that the formally
preferred model also provides a physically plausible description of the VSF
shape across the fitted radial range.
